% Part: lambda-calculus
% Chapter: church-rosser
% Section: parallel-beta-reduction

\documentclass[../../../include/open-logic-section]{subfiles}

\begin{document}

\olfileid[pt]{lam}{cr}{pb}

\olsection{$\beta$-redução paralela}

Introduzimos a noção de \emph{$\beta$-redução paralela} e provamos que ela
tem a propriedade de Church--Rosser.

\begin{defn}[$\beta$-redução paralela, $\bredpar$] \ollabel{defn:bredpar}
  A redução paralela ($\bredpar$) de termos é definida indutivamente da
  seguinte maneira:
  \begin{enumerate}
    \item \ollabel{defn:bredpar1} $x \bredpar x$.
    \item \ollabel{defn:bredpar2} Se $N \xrightarrow{\beta} N'$, então
      $\lambd[x][N] \bredpar \lambd[x][N']$.
    \item \ollabel{defn:bredpar3} Se $P \bredpar P'$ e $Q \bredpar Q'$, então
      $PQ \bredpar P'Q'$.
    \item \ollabel{defn:bredpar4} Se $N \bredpar N'$ e $Q \bredpar Q'$, então
      $(\lambd[x][N])Q \bredpar \Subst{N'}{Q'}{x}$.
  \end{enumerate}
\end{defn}

A $\beta$-redução paralela nos permite reduzir qualquer número de redexes
em um termo em um único passo. Ela difere da $\beta$-redução porque só
podemos contrair redexes presentes no termo original, e não redexes que
surjam da própria $\beta$-redução paralela. Por exemplo, o termo
$(\lambd[f][fx])(\lambd[y][y])$ só pode ser $\beta$-reduzido paralelamente
a si próprio ou a $(\lambd[y][y])x$, mas não mais adiante até~$x$, embora
ele se $\beta$-reduza a~$x$, pois esse redex só surge depois de um passo de
$\beta$-redução paralela. Um segundo passo de $\beta$-redução paralela,
contudo, produz~$x$.

\begin{thm}\ollabel{thm:refl}
  $M \bredpar M$.
\end{thm}

\begin{proof}
  Exercício.
\end{proof}

\begin{prob}
Demonstre \olref[lam][cr][pb]{thm:refl}.
\end{prob}

\begin{defn}[desenvolvimento $\beta$-completo]\ollabel{defn:bcd}
  O \emph{desenvolvimento $\beta$-completo}~$\bcd{M}$ de $M$ é definido
  indutivamente da seguinte maneira:
  \begin{align}
    \bcd{x} &= x \ollabel{defn:bcd1} \\
    \bcd{(\lambd[x][N])} &= \lambd[x][\bcd{N}] \ollabel{defn:bcd2}\\
    \bcd{(PQ)} &= \bcd{P}\bcd{Q} && \text{se $P$ não é uma abstração-$\lambd$}
    \ollabel{defn:bcd3} \\
    \bcd{((\lambd[x][N])Q)} &= \Subst{\bcd{N}}{\bcd{Q}}{x} \ollabel{defn:bcd4}
  \end{align}
\end{defn}

O desenvolvimento $\beta$-completo de um termo é, como sugere o nome, uma
``redução paralela completa''. Enquanto na $\beta$-redução paralela ainda
podemos optar por não contrair um redex, no desenvolvimento completo somos
obrigados a contrair todos eles. Assim, o desenvolvimento completo de
$(\lambd[f][fx])(\lambd[y][y])$ é $(\lambd[y][y])x$, e não o próprio termo.

\begin{editorial}
  Esta definição tem o problema de que ainda não introduzimos como definir
  recursivamente funções sobre termos (do $\lambd$). Isso será corrigido no
  futuro.
\end{editorial}

\begin{lem}\ollabel{lem:comp}
  Se $M \bredpar M'$ e $R \bredpar R'$, então $\Subst{M}{R}{y}
  \bredpar \Subst{M'}{R'}{y}$.
\end{lem}

\begin{proof}
  Por indução na !!{derivation} de $M \bredpar M'$.
  \begin{enumerate}
    \item O último passo é \olref{defn:bredpar1}: exercício.
    \item O último passo é \olref{defn:bredpar2}: então $M$ é
      $\lambd[x][N]$ e $M'$ é $\lambd[x][N']$,
      onde $N \bredpar N'$. Queremos provar que
      $\Subst{(\lambd[x][N])}{R}{y} \bredpar
      \Subst{(\lambd[x][N'])}{R'}{y}$, isto é,
      $\lambd[x][\Subst{N}{R}{y}] \bredpar
      \lambd[x][\Subst{N'}{R}{y}]$. Isso decorre imediatamente de
      \olref{defn:bredpar2} e da hipótese de indução.
    \item O último passo é \olref{defn:bredpar3}: exercício.
    \item O último passo é \olref{defn:bredpar4}: $M$ é
      $(\lambd[x][N])Q$ e $M'$ é $\Subst{N'}{Q'}{x}$. Queremos provar que
      $\Subst{((\lambd[x][N])Q)}{R}{y} \bredpar
      \Subst{\Subst{N'}{Q'}{x}}{R'}{y}$, isto é,
      $(\lambd[x][\Subst{N}{R}{y}])\Subst{Q}{R}{y} \bredpar
      \Subst{\Subst{N'}{R'}{y}}{\Subst{Q'}{R'}{y}}{x}$. Isso decorre de
      \olref{defn:bredpar4} e da hipótese de indução.
  \end{enumerate}
\end{proof}

\begin{prob}
  Complete a prova de \olref[lam][cr][pb]{lem:comp}.
\end{prob}

\begin{lem}\ollabel{lem:cont}
  Se $M \bredpar M'$, então $M' \bredpar \bcd{M}$.
\end{lem}
\begin{proof}
  Por indução na !!{derivation} de $M \bredpar M'$.
  \begin{enumerate}
    \item A última regra é \olref{defn:bredpar1}: exercício.
    \item A última regra é \olref{defn:bredpar2}: $M$ é
      $\lambd[x][N]$ e $M'$ é $\lambd[x][N']$, com
      $N \bredpar N'$. Queremos mostrar que $\lambd[x][N'] \bredpar
      \bcd{(\lambd[x][N])}$, isto é, $\lambd[x][N'] \bredpar
      \lambd[x][\bcd{N}]$ por \olref{defn:bcd2}. Isso decorre de
      \olref{defn:bredpar2} e da hipótese de indução.
    \item A última regra é \olref{defn:bredpar3}: $M$ é $PQ$ e $M'$ é
      $P'Q'$ para alguns $P$, $Q$, $P'$ e $Q'$, com $P \bredpar P'$ e
      $Q \bredpar Q'$. Pela hipótese de indução, temos $P' \bredpar \bcd{P}$
      e $Q' \bredpar \bcd{Q}$.
      \begin{enumerate}
        \item Se $P$ é $\lambd[x][N]$ para alguns $x$ e $N$, então $P'$ deve
          ser $\lambd[x][N']$ para algum $N'$ com $N \bredpar N'$. Pela
          hipótese de indução, temos $N' \bredpar \bcd{N}$ e
          $Q' \bredpar \bcd{Q}$. Então $(\lambd[x][N'])Q' \bredpar
          \Subst{\bcd{N}}{\bcd{Q}}{x}$ por \olref{defn:bredpar4}.
        \item Se $P$ não é uma abstração-$\lambd$, então
          $P'Q' \bredpar \bcd{P}\bcd{Q}$ por \olref{defn:bredpar3}, e o
          lado direito é $\bcd{PQ}$ por \olref{defn:bcd3}.
      \end{enumerate}
    \item A última regra é \olref{defn:bredpar4}: $M$ é
      $(\lambd[x][N])Q$ e $M'$ é $\Subst{N'}{Q'}{x}$ para alguns $x$,
      $N$, $Q$, $N'$ e~$Q'$, com $N \bredpar N'$ e $Q \bredpar Q'$. Pela
      hipótese de indução, sabemos que $N' \bredpar \bcd{N}$ e
      $Q' \bredpar \bcd{Q}$. Por \olref{lem:comp}, temos
      $\Subst{N'}{Q'}{x} \bredpar \Subst{\bcd{N}}{\bcd{Q}}{x}$, cujo lado
      direito é exatamente $\bcd{((\lambd[x][N])Q)}$.
  \end{enumerate}
\end{proof}

\begin{prob}
  Complete a prova de \olref[lam][cr][pb]{lem:cont}.
\end{prob}

\begin{thm}\ollabel{thm:cr}
  $\bredpar$ tem a propriedade de Church--Rosser.
\end{thm}

\begin{proof}
  Imediato de \olref{lem:cont}.
\end{proof}

\end{document}
